\section{谱偏差怎样影响 PCA 与随机投影}
\label{sec:11-Mathematical-Statistics/09-Random-Matrix-Theory/03}

同一份高维数据摆在面前，两件事同时在发生。你用 PCA 取前几个主方向，得到的方向在重采样时来回摆动，报出去的``解释方差比''还偏高；你用一个随机高斯矩阵把同一份数据压到几百维做近邻检索，召回率却几乎没掉。一个是精心估计出来的方向不可靠，一个是随手抽出来的方向很好用，直觉上该反过来才对。

上一节给了噪声的形状和信号的门槛。这一节回答门槛之后的事，顺便解释这两件事为什么不矛盾。

\subsection{信号强度跨过 \(\sqrt c\) 之前，主方向估计与真方向渐近正交}

把问题写成最简单的形式：真协方差是单位阵加一个方向上的凸起，

\[
\Sigma=I+h\,vv^{\trans},\qquad \norm{v}=1,\ h>0 .
\]

只有 \(v\) 这一个真实主方向，强度由 \(h\) 控制，其余方向全是各向同性噪声。样本量与维数之比仍记 \(c=p/n\)。记 \(\widehat v\) 是样本协方差最大特征值对应的特征向量，问它离 \(v\) 有多远。

答案由 Baik--Ben Arous--Péché 相变给出，分两段：

\[
\begin{aligned}
h\le\sqrt c:\quad &\widehat\lambda_{\max}\to(1+\sqrt c)^2, &\quad \abs{\ip{\widehat v}{v}}^2&\to 0,\\
h>\sqrt c:\quad &\widehat\lambda_{\max}\to(1+h)\Bigl(1+\frac ch\Bigr), &\quad \abs{\ip{\widehat v}{v}}^2&\to\frac{1-c/h^2}{1+c/h}.
\end{aligned}
\]

上半段是这一节最该记住的一行。信号强度不到 \(\sqrt c\) 时，最大特征值停在 MP 上界上，与纯噪声毫无区别；而估计出的主方向与真方向的内积趋于零——在高维里，两个随机方向本来就近乎正交，\(\widehat v\) 退化成了这样一个方向。它不是``估得不太准''，它是与真方向无关。信号确实存在，样本量确实在增长，只要 \(p/n\) 不降，你就永远拿不到它。

取 \(c=0.5\) 感受一下数值。阈值是 \(\sqrt{0.5}=0.707\)。\(h=0.7\) 时方向完全丢失；\(h=0.75\) 时 \(\abs{\ip{\widehat v}{v}}^2=0.067\)，夹角 \(75^\circ\)；\(h=1\) 时夹角 \(55^\circ\)；\(h=1.5\) 时仍有 \(40^\circ\)；要把夹角压到 \(22^\circ\) 需要 \(h=4\)，也就是信号方向的方差是噪声方向的五倍。

\begin{center}
\begin{tikzpicture}[>=stealth,x=1.5cm,y=0.038cm]
\fill[black!10] (0,0) rectangle (0.707,90);
\draw[->] (-0.1,0) -- (4.5,0) node[right,font=\small] {\(h\)};
\draw[->] (0,-4) -- (0,102) node[above,font=\small] {\(\angle(\widehat v,v)\)};
\foreach \y in {0,30,60,90}
  \draw (0,\y) -- (-0.07,\y) node[left,font=\scriptsize] {\(\y^\circ\)};
\foreach \x/\l in {1/1,2/2,3/3,4/4}
  \draw (\x,0) -- (\x,-3) node[below,font=\scriptsize] {\l};
\draw[dashed] (0.707,0) -- (0.707,98);
\node[font=\scriptsize,anchor=south] at (0.707,98) {\(h=\sqrt c\)};
\draw[very thick] (0,90) -- (0.707,90);
\draw[very thick,smooth] plot coordinates {(0.707,90) (0.75,75) (0.8,68.5) (0.9,60.3)
  (1.0,54.7) (1.2,47.3) (1.5,40.2) (2.0,33.2) (2.5,28.9) (3.0,25.9) (4.0,21.9)};
\node[font=\scriptsize,anchor=west] at (1.55,84) {\(h\le\sqrt c\)：与真方向渐近正交};
\draw[->] (1.5,83) -- (0.45,88);
\node[font=\scriptsize,anchor=west] at (2.35,55) {越过阈值后才开始};
\node[font=\scriptsize,anchor=west] at (2.35,46) {逐渐对准，但很慢};
\end{tikzpicture}
\end{center}

\emph{\(c=0.5\) 时，估计主方向与真方向的夹角在 \(h=\sqrt c\) 处从 \(90^\circ\) 折下来，此后衰减得相当缓慢。}

图里那个折点是硬的：左边是一整段平台，方向信息为零；右边虽然开始下降，但曲线一路贴得很高。碎石图上一个刚刚越过 MP 上界的特征值，对应的主方向可能与真方向差着 \(70^\circ\) 以上——特征值已经``显著''，特征向量还什么都不是。这也说明上一节那条阈值线该怎么读：它筛的是``这里可能有信号''，不是``这个方向可以拿去用''。

\subsection{越过阈值之后，特征值仍然偏高，但偏高的量可以算回去}

相变公式的下半段还带来一个可用的副产品。真实的最大特征值是 \(1+h\)，而样本给出的是 \((1+h)(1+c/h)\)，比值 \(1+c/h\) 就是高估的倍数。\(c=0.5\)、\(h=1.5\) 时高估 \(33\%\)；\(c=1\) 时同一个 \(h\) 高估 \(67\%\)。信号越弱、维数比越大，虚高越严重——而弱信号恰恰是最容易被过度解读的那一类。

好在这个偏差是可逆的。把 \(\ell=(1+h)(1+c/h)\) 当成关于 \(h\) 的方程整理，得到 \(h^2+(1+c-\ell)h+c=0\)，于是

\[
\widehat h=\frac{(\ell-1-c)+\sqrt{(\ell-1-c)^2-4c}}{2},
\]

判别式非负当且仅当 \(\ell\ge(1+\sqrt c)^2\)，与相变阈值自洽。观测到 \(\ell=3.33\)、\(c=0.5\)，回代得 \(\widehat h=1.5\)，真实的最大特征值是 \(2.5\) 而不是 \(3.33\)。同一个思路可以用到解释方差比上：先把每个越界特征值去偏，再算比例，报出去的数字才有意义。用原始的样本特征值算比例，等于把噪声吸进去的那部分方差也记在主成分账上，本单元第一节开头那份基因数据里``前三个主成分解释了 \(80\%\) 方差''的说法，虚高多半就是这么来的。

去偏能修的只有特征值。特征向量的夹角由 \(\abs{\ip{\widehat v}{v}}^2\) 那个公式给出，它是信息量的上限，没有哪种后处理能凭空补上——除非引入额外结构，比如假设 \(v\) 稀疏，用稀疏 PCA 把搜索限制在少数坐标上，把有效维数降下来。这一点与 \hyperref[sec:03-Linear-Algebra/07-SVD-and-Data-Applications/03]{``低秩近似：用少数方向保留主要结构''} 里``奇异值稳定、奇异向量未必稳定''是同一件事的两种来源：那里是谱隙小导致的病态，这里是随机性导致的失真。

\subsection{随机方向不需要估计，所以不被 \(p/n\) 拖累}

现在换一件事做。不找主方向，直接取一个 \(k\times p\) 的随机矩阵 \(R\)，元素独立同分布、方差 \(1/k\)，把每个点映成 \(Rx\)。Johnson--Lindenstrauss 引理说：给定 \(n\) 个点和精度 \(\varepsilon\in(0,1)\)，只要

\[
k\ \ge\ \frac{4\ln n}{\varepsilon^2/2-\varepsilon^3/3},
\]

就以高概率对所有 \(\binom n2\) 个点对同时满足 \((1-\varepsilon)\norm{x_i-x_j}^2\le\norm{Rx_i-Rx_j}^2\le(1+\varepsilon)\norm{x_i-x_j}^2\)。

证明骨架只有两步。对固定的一个向量 \(u\)，\(\norm{Ru}^2/\norm{u}^2\) 是 \(k\) 个独立平方项的平均，期望为 1，偏离 \(\varepsilon\) 的概率按 \(\mathrm e^{-k\varepsilon^2/8}\) 量级衰减；然后对 \(\binom n2\) 个差向量做联合界，只要 \(k\) 与 \(\ln n\) 成正比，失败概率就压得住。条件对账里最值得注意的是没出现什么：没有对数据分布的假设，没有对 \(p\) 的限制，也没有任何需要从数据里估计的量。

\(k\) 与 \(p\) 无关，这是关键的一句。一份 \(p=100000\) 维的 TF--IDF 表示、\(n=10^5\) 篇文档，取 \(\varepsilon=0.3\)，公式给出 \(k\approx1280\)。维数降了近八十倍，而任意两篇文档之间的距离误差不超过三成。把 \(p\) 换成一百万，\(k\) 一个不变。实践中理论常数偏保守，几百维往往就够，这也是近邻检索里先随机投影再建索引的常规做法：投影后的低维空间上跑 KD 树或哈希桶，候选集再回到原空间精排，召回损失可控。随机化 SVD 也是同一个动作的另一种用法，先用随机投影抓住列空间，再在小矩阵上做精确分解，\hyperref[sec:08-Numerical-Analysis/05-Numerical-Linear-Algebra/06]{``SVD、低秩与数值秩''}讲了它的误差账。

把两件事并排放，落差就清楚了。PCA 要从数据里估计方向，估计的质量受 \(p/n\) 支配，\(c\) 一大就退化到与真方向正交；随机投影不估计任何东西，它的保证是随机矩阵自身的集中性质，与数据、与 \(p\) 都无关，让出的只是一个可以事先选定的失真 \(\varepsilon\)。用随机矩阵谱的语言说，这是同一个现象的两侧：一个 \(k\times d\) 的随机矩阵，当比值 \(d/k\) 小的时候，它的奇异值全都挤在 1 附近，几乎是等距的。PCA 的处境是 \(c=p/n\) 大，谱被拉开；随机投影则是自己把比值调小——只要 \(k\) 远大于需要保住的那点信息量（\(\ln n\) 级），比值就小，谱就贴着 1。高维不是一律有害，害处集中在``需要用有限样本估计一个高维方向''这件事上。

\subsection{两条结论合起来给出一条分工}

于是有了一条可操作的分工。任务如果要的是可解释的方向、要报``第一主成分代表什么''，那它对谱失真敏感，先算 \(c\)，用 MP 上界数一数有几个特征值越界，对越界的特征值去偏，再看去偏后的夹角公式允许你相信到什么程度；\(c\) 大而信号弱时，诚实的结论是这个方向拿不出来。任务如果要的只是距离、近邻、聚类这类几何量，方向本身不用解释，那就别去估计方向，随机投影给的保证更硬，而且不随 \(p\) 恶化。协方差收缩介于两者之间：它承认样本谱不可信，用一个与 \(c\) 挂钩的强度把谱往回压，\hyperref[sec:13-Machine-Learning/06-Probabilistic-Models/02]{``LDA 与收缩：判别方向的几何与高维失效''}在分类任务上算得更细。

有一处容易越界的推论要拦一下。随机投影压到几百维之后模型表现不掉，不能反推``数据的内在维度只有几百''——JL 引理对任意 \(n\) 个点都成立，不需要它们躺在任何低维流形上，\hyperref[sec:14-Deep-Learning/10-Interpretability-Mathematics/03]{``高维数据的低维结构检验''}把这个混淆拆开过。压缩成功只说明当前任务对这次压缩稳健，不说明数据有低维结构。

本单元到此把一件事讲完了：\(p\) 与 \(n\) 可比时，样本协方差的谱有确定的失真形状，失真可以计算、部分可以校正，剩下不可校正的部分则界定了高维推断的边界。同一套随机矩阵的谱工具还在别处出现——核矩阵的谱、随机特征映射的近似误差、深层网络的 Jacobian 奇异值分布，用的都是这里的语言。
